\Subhead{4.7.1}{Резонаторный усилитель и усилитель бегущей волны. Полоса пропускания усилителя бегущей волны}

В резонаторном усилителе электронный пучок или активный элемент взаимодействует с полем одной/нескольких резонансных мод. Энергия пучка превращается в СВЧ-мощность, когда ток имеет компоненту нужной фазы относительно поля: средняя работа поля над электронами отрицательна. Резонатор накапливает поле, но ограничивает полосу примерно $\Delta\omega\sim\omega_0/Q_L$ в простейшей одно-модовой модели. Многорезонаторная структура может расширить полосу ценой настройки и других компромиссов.

В усилителе бегущей волны обмен распределён по длине. Малое продольное поле модулирует скорость электронов; группировка создаёт переменный ток, который усиливает синхронную волну. Условие $\omega-k_zv_0\simeq0$ удерживает электронные сгустки в тормозящей фазе. Входной сигнал задаёт частоту и начальную модуляцию; энергия усиления поступает от постоянного ускоряющего напряжения. Предельная мощность связана с $I_0V_0$ и КПД, а не с энергией входной волны.

Оценка полосы следует из фазового рассогласования:
\begin{equation}
 \Delta\phi(\omega)=[\beta_c(\omega)-\omega/v_0]L.
\end{equation}
Если $|\Delta\phi|$ становится большим, вклад разных участков перестаёт складываться. Около синхронизма
\begin{equation}
 \frac{\dd\Delta\phi}{\dd\omega}=L\left(\frac1{v_{g,c}}-\frac1{v_0}\right).
\end{equation}
Поэтому малая дисперсия замедляющей структуры и согласование фазовой скорости в широкой полосе дают широкополосность. Это оценка когерентного взаимодействия, а не точная формула полосы высокоусилительной ЛБВ: связь с пучком, пространственный заряд, потери и отражения меняют дисперсионные корни. Самовозбуждение при паразитной обратной связи ограничивают развязкой и поглотителем, не нарушая прямое взаимодействие.
